\section{Limitations}
\label{sec:limits}

\textbf{Scale.} The largest model trained is \largeParams{}M parameters, with five strategies at two migration points;
the full sweep stops at \paramsM{}M. Everything ran on one laptop GPU and one small text dataset. The
mechanisms of Section~\ref{sec:transient} depend only on the decay rates and are not specific to a size, and the
measured picture is stable across a factor of a hundred in parameter count (Fig.~\ref{fig:frontier},
Section~\ref{sec:large}), but none of this
is evidence about a multi-billion-parameter run. Every figure quoted for such a model is an extrapolation and is
labelled.

\textbf{Three seeds.} Each cell has \nSeedsMax{} seeds, so each seed band is the 95th percentile over three pairs of
controls and each interval resamples three units. The three groups of strategies are separated by more than an order
of magnitude, which three seeds resolve; the differences inside a group, such as rank one against rank four, are not
resolved and we do not claim them. The dependence on $\beta_2$ was measured at one model size and two migration
points, and $\beta_1$ was not varied. The vocabulary is a few dozen characters, so the embedding results do not transfer to a
tokenised model.

\textbf{Simulated transport.} A migration here serialises to local storage, tears the model and optimizer down and
rebuilds them from the bundle. That reproduces exactly the information lost by a real migration, which is what the
loss measurements depend on. It does not reproduce a network: the drain times in Section~\ref{sec:policy} are bytes
divided by an assumed bandwidth, and the measured latencies of Section~\ref{sec:latency} cover encoding and local I/O
only. Section~\ref{sec:link} sends payloads between two processes through a paced socket, which reproduces a
bandwidth and not a network: there is no loss, jitter or competing traffic. Payloads were also resumed on a second
machine (Section~\ref{sec:device}), but that transfer was not timed. A timed migration between two machines over a real
link has not been run.

\textbf{One host.} The multi-rank runtime was exercised with CPU ranks on one machine over Gloo. Resizing a job changes
its world size, not its node count, and no claim is made about a GPU cluster.

\textbf{Device classes.} Four trainable devices were tested: an x86 CPU, an NVIDIA laptop GPU, an Intel integrated GPU
and an NVIDIA T4 on a second machine.
The laptop also has a neural processing unit, and it was not used: current NPU software stacks execute inference
graphs and offer no backward pass~\cite{openvino2025npu,feng2024steppingforward}, so a training job cannot be migrated
to one. No TPU run was made. The integrated-GPU and remote-GPU runs use other framework releases than the
rest, so those rows change device and framework version together.

\textbf{One optimizer family.} All runs use AdamW. The analysis carries over to any method built on the same two
exponential averages; it says nothing about second-order or sign-based methods, or about schedules with restarts.

\textbf{Sharded state.} Rank zero drains the whole payload. With tensor- or pipeline-parallel training the state is
partitioned across devices and each shard would drain its own part; the strategies apply per tensor and do not depend
on where a tensor lives, but that path was not built.

\textbf{Steps lost.} The cost model's $\lambda_s$ comes from training loss and has a wide noise band. The peak, regret
and recovery measures on held-out loss are the reliable ones, and our conclusions rest on those.

\textbf{Concurrent work.} The literature search was last run in November 2026. The two closest preprints are discussed
in Section~\ref{sec:related}; others may have appeared since.
